\section{GaitNet}
\label{sec:methodology-gaitnet}

GaitNet is designed to isolate the optimal footstep action from a
continuous configuration space using gradient-based optimization.
Instead of selecting from a discrete grid populated by a prior evaluation
network, GaitNet serves as the primary reward estimator, enabling direct
optimization in the continuous output space.

Specifically, GaitNet optimizes the one-hot encoded leg index
$\mathbf{f_c}^{\prime}$
and the spatial coordinates via projected gradient ascent. The inclusion of a
no-action candidate, evaluated directly alongside the optimized paths,
allows the greedy generation of acyclic gaits in which multiple feet
can be in the swing phase simultaneously.

%%%%%%%%%%%%%%%%%%%%%%%%%%%%%%%%%%%%%%%%%%%%%%%%%%%%%%%%%%%%%%%%%%%%%%%%%%%%%%%%
\subsection{Architecture}
\label{subsec:methodology-gaitnet-architecture}

GaitNet is responsible for ranking footstep candidates based on the
one-hot encoded leg index $\mathbf{f_c}^{\prime}$ and the GaitNet
input $\mathbf{x_g}$:

\begin{equation}
  \mathbf{x_g} =
  \begin{bmatrix}
    \mathbf p_{b,xy} \\
    \mathbf r_{w,z} \\
    \mathbf v_b \\
    \mathbf \omega_b \\
    \mathbf u \\
    \mathbf g \\
    \mathbf c
  \end{bmatrix}
  \label{eq:gaitnet-x}
\end{equation}

Descriptions of the first five terms dictate the primary robot state,
tracking current
base position $p_{b,xy}$ and rotation $r_{w,z}$ in world coordinates,
body linear
$\mathbf v_b$ and angular $\mathbf \omega_b$ velocities, and current
foot positions $\mathbf u$.
The additional terms are $\mathbf g$ and $\mathbf c$, being the
gravity vector and contact state
respectively. Using $\mathbf{x_g}$ and $\mathbf{f_c}^{\prime}$,
GaitNet outputs a logit associated with
each evaluated spatial coordinate and the desired swing duration if
that step is selected.

\begin{figure}[H]
  \centering
  \includegraphics[width=0.75\linewidth]{images/diagrams/gait-network-architecture.png}
  \caption{GaitNet architecture. Sections of the diagram include the
    footstep candidate encoder in red, robot state encoder in blue, and
    shared trunk in green. The final output is a logit encoding the
  value of this option and the desired swing duration if this action is taken.}
  \label{fig:diagram-gaitnet-architecture}
\end{figure}

The model is designed to jointly evaluate robot state and footstep
candidates with a middle fusion architecture \cite{feng2021deep}. The
architecture (shown in \autoref{fig:diagram-gaitnet-architecture})
consists of two encoders, a shared trunk, and two task-specific output heads:

\begin{itemize}
  \item Robot state encoder: The robot state vector is processed by a
    three-layer feedforward network with intermediate dimensionality
    of 128 and ReLU activations. This produces a fixed-dimensional
    latent representation of the current robot state.
  \item Footstep encoder: Each footstep candidate is encoded by a
    two-layer feedforward network with hidden size 64 and ReLU
    activations. No-action candidates are represented through a fixed embedding.
  \item Shared trunk: The concatenated robot state and footstep
    embeddings are processed by a three-layer feedforward trunk,
    reducing dimensionality while applying layer normalization and
    ReLU nonlinearity.
  \item Output heads: Two parallel prediction heads are applied to
    the shared trunk:
    \begin{itemize}
      \item Logit head: A single-layer network outputs the predicted
        reward value as a logit.
      \item Duration head: A single-layer network with sigmoid
        activation outputs a normalized swing duration, which is then
        scaled to the range (0.1, 0.3)\,s.
    \end{itemize}
\end{itemize}

This design allows the network to assign an expected value and a feasible
swing duration to any given spatial coordinate, while explicitly supporting a
no-action option via a dedicated embedding.

%%%%%%%%%%%%%%%%%%%%%%%%%%%%%%%%%%%%%%%%%%%%%%%%%%%%%%%%%%%%%%%%%%%%%%%%%%%%%%%%
\subsection{Projected Gradient Ascent Optimization}
\label{subsec:methodology-gaitnet-gradient-ascent}

Instead of sampling from a discrete pool of generated candidates, we formulate
the footstep selection as a continuous optimization problem over the spatial
coordinates of the footstep action. GaitNet's reward prediction logit is treated
as a differentiable objective function maximized via projected gradient ascent.

For each foot, $N=4$ random initial spatial samples are drawn uniformly from the
reachable kinoynamic workspace. For each sample, the spatial
coordinates are updated
in parallel for 4 steps to maximize the expected value output logit
computed by GaitNet.
After each gradient update step, the new spatial coordinates are
projected back to the
nearest valid configuration within the reachable stepping continuous
space (cspace)
to ensure physical validity. This forces the generated footsteps to
avoid un-steppable
terrain like holes while preserving maximum reward.

Let the fixed "no-action" embedding be evaluated alongside the spatial updates.
After the 4 projection and ascent steps, the candidate with the highest reward
logit across all $4 \times 4$ options (4 tracks, 4 legs) is compared to the
predicted value of the "no-action" candidate. The best resulting footstep or
no-action command is then deployed to the inner controller. This
optimization runs
at control frequencies, tightly integrating flexible learning-driven evaluation
with robust multi-leg acyclic gait coordination.

%%%%%%%%%%%%%%%%%%%%%%%%%%%%%%%%%%%%%%%%%%%%%%%%%%%%%%%%%%%%%%%%%%%%%%%%%%%%%%%%
\subsection{Training}
\label{subsec:methodology-gaitnet-training}

Due to the inclusion of the $\mathbf{c}$ term in the input, GaitNet
cannot be effectively trained using supervised learning because of
the high dimensionality of the problem. Instead, it is trained using
PPO \cite{Schulman_proximal_2017}. In this setup, the actor network
is GaitNet (\autoref{fig:diagram-gaitnet-architecture}), while the
critic follows a similar architecture but omits the duration head.
All logits from the critic are passed through a final MLP to produce
a single value estimate. This MLP consists of two hidden layers of
size 64, with ReLU activations applied to all layers except the output.

A custom actor-critic implementation was necessary to handle
GaitNet's dual outputs. The logits define a categorical distribution
over footstep candidates, while the duration output is treated as a
separate normal distribution with a fixed standard deviation of
0.01\,s for each action. To prevent multiple no-action candidates
from skewing the selection, all but one no-action candidate are
removed (set to $-\infty$).

The total log probability of an action is computed as the sum of the
categorical log probability of the selected footstep candidate and
the log probability of the duration under the normal distribution.
During action sampling, a footstep candidate is first sampled from
the categorical distribution, followed by sampling the duration from
its associated normal distribution. The footstep candidate and
duration are then combined to form a complete footstep action.

Following the reward structure, termination conditions, commands, and
hyperparameters outlined in
\autoref{chap:appendix-gaitnet-training-configuration}, GaitNet is
trained for 700 environment steps using 100 agents in parallel.
Each environment step simulates 20\,s. The 700 step end point was chosen
because improvement drastically slowed by then.
